\documentclass[11pt]{article}

\usepackage[a4paper, margin=28mm]{geometry}
\usepackage{amsmath, amssymb}
\usepackage{booktabs}
\usepackage{xcolor}
\usepackage{listings}
\usepackage{microtype}
\usepackage[hidelinks]{hyperref}

\definecolor{codegray}{gray}{0.45}
\definecolor{rulegray}{gray}{0.75}

\lstdefinelanguage{Rust}{
  keywords={fn, let, mut, if, else, while, for, in, match, return, struct,
    enum, impl, trait, pub, use, const, loop, continue, break, ref, type,
    where, as, self, Self, true, false, Some, None, Ok, Err},
  sensitive=true,
  comment=[l]{//},
  string=[b]",
}
\lstset{
  language=Rust,
  basicstyle=\small\ttfamily,
  commentstyle=\color{codegray}\itshape,
  keywordstyle=\bfseries,
  columns=fullflexible,
  keepspaces=true,
  breaklines=true,
  frame=leftline,
  rulecolor=\color{rulegray},
  framesep=8pt,
  xleftmargin=12pt,
  aboveskip=10pt,
  belowskip=10pt,
}

% loc counter shown under each section heading: \loc{section name}{lines}{cumulative}
\newcommand{\loc}[3]{\noindent{\small\color{codegray}\texttt{#1}: #2 lines, running total #3 of 999}\par\medskip}

\title{\textbf{nanospice}\\[4pt]\large A SPICE circuit simulator in 1000 lines of Rust}
\author{Milan Rother}
\date{2026}

\begin{document}
\maketitle

\begin{abstract}
\noindent
nanospice is a circuit simulator with the classic Berkeley SPICE algorithm
set: modified nodal analysis, Newton-Raphson with junction limiting and two
operating point fallbacks, trapezoidal transient integration with local
truncation error timestep control, and small-signal AC analysis. The entire
solver lives in one Rust source file of at most 1000 nonblank, noncomment
lines, enforced by a test, with no dependencies outside the standard
library. This report derives each algorithm, explains the design decisions
the budget forced, and maps both to the code. Each section header states how
many of the 999 lines the corresponding code section spends.
\end{abstract}

\section{The budget as a design tool}

The rule: one file, \texttt{src/main.rs}, at most 1000 lines after removing
blank lines and comments, standard library only. A test reads the file,
counts, and fails the build above the limit. Tests, example netlists, and
comments are free. The point of the rule is not minimalism for its own sake.
A hard budget forces every feature to displace another, which makes the cost
of each decision explicit, and it keeps the whole simulator readable in one
sitting, which is the educational goal of the repository.

Table~\ref{tab:budget} shows where the lines go. The parser is the largest
consumer, which is typical: the numerics of a circuit simulator are compact,
the input format is not.

\begin{table}[h]
\centering
\begin{tabular}{lrr}
\toprule
section & lines & cumulative \\
\midrule
constants, error and output plumbing & 23 & 23 \\
numbers: \texttt{Num} trait, complex, unit suffixes & 78 & 101 \\
circuit data model & 63 & 164 \\
parser & 262 & 426 \\
sparse LU & 96 & 522 \\
device models & 50 & 572 \\
solver: assembly, Newton, operating point & 179 & 751 \\
output selection & 31 & 782 \\
analyses and main & 217 & 999 \\
\bottomrule
\end{tabular}
\caption{Line budget by code section, in file order.}
\label{tab:budget}
\end{table}

Everything the budget excluded is listed with reasons in
section~\ref{sec:cuts}.

\section{Modified nodal analysis}
\loc{solver, shared with sections 4 and 5}{179}{751}

The unknown vector stacks the node voltages and one branch current for every
voltage-defined element:
\[
x = (v_1, \dots, v_N,\; i_1, \dots, i_B)^T ,
\]
where the branches belong to voltage sources, inductors, and VCVS elements.
Kirchhoff's current law gives one equation per node; each voltage-defined
element adds one equation of the form $v_p - v_n = E$ and contributes its
current $\pm i$ to the node equations. A two-terminal conductance $g$
between nodes $p$ and $n$ stamps the familiar quad
\[
\begin{pmatrix} +g & -g \\ -g & +g \end{pmatrix}
\]
into rows and columns $(p, n)$. A VCCS is the same quad with the rows taken
at the output nodes and the columns at the controlling nodes, so one helper
covers both:

\begin{lstlisting}
fn quad(&mut self, p: usize, n: usize, cp: usize, cn: usize, g: T) {
    self.add(p, cp, g);
    self.add(p, cn, T::zero() - g);
    self.add(n, cp, T::zero() - g);
    self.add(n, cn, g);
}
fn branch(&mut self, p: usize, n: usize, br: usize) {
    self.add(p, br, T::one());
    self.add(n, br, T::zero() - T::one());
    self.add(br, p, T::one());
    self.add(br, n, T::zero() - T::one());
}
\end{lstlisting}

\paragraph{The ground row.}
Ground is unknown index 0, and it is a dummy: \texttt{add} silently drops
any entry with a row or column index of 0, and the solver loops start at
index 1 with $x_0 = 0$ pinned. The alternative is the usual
\emph{if node is not ground} test inside every stamp, roughly two dozen
branches across the device set. Spending one guard inside \texttt{add}
instead of one per stamp is the kind of trade the budget rewards, and it
also removes a class of copy-paste bugs.

\section{Sparse LU}
\loc{sparse LU}{96}{522}

Dense LU with partial pivoting is about 25 lines and was the first
implementation. It is $O(n^3)$, which caps practical circuits at a few
hundred nodes. The replacement is a genuinely sparse LU in 96 lines,
including the generic data structure.

Each matrix row is a vector of $(\text{column}, \text{value})$ pairs kept
sorted by column. The factorization maintains one invariant:

\begin{quote}
Before elimination step $k$, every row stored at index $i \ge k$ contains
only columns $\ge k$.
\end{quote}

The invariant holds initially because column 0 is never stored, and step
$k$ preserves it because eliminating column $k$ from a row removes that
entry. Its consequence is the whole design: if a row at index $i \ge k$ has
an entry in column $k$ at all, that entry is the \emph{first} element of its
vector. Pivot search reduces to scanning first elements, and the update of a
row against the pivot row is a merge of two sorted sequences in which
fill-in appears automatically as the union of the column sets:

\begin{lstlisting}
// row a minus f times pivot row p, both sorted
fn merge<T: Num>(a: &[(usize, T)], p: &[(usize, T)], f: T) -> Vec<(usize, T)> {
    let mut o = Vec::with_capacity(a.len() + p.len());
    let (mut i, mut j) = (0, 0);
    while i < a.len() && j < p.len() {
        if a[i].0 == p[j].0 {
            o.push((a[i].0, a[i].1 - f * p[j].1));
            i += 1;
            j += 1;
        } else if a[i].0 < p[j].0 {
            o.push(a[i]);
            i += 1;
        } else {
            o.push((p[j].0, T::zero() - f * p[j].1));
            j += 1;
        }
    }
    o.extend_from_slice(&a[i..]);
    for &(c, v) in &p[j..] {
        o.push((c, T::zero() - f * v));
    }
    o
}
\end{lstlisting}

Pivoting is by largest magnitude among the candidate first entries, with
rows swapped whole. MNA needs this: the constraint row of a voltage source
has a structural zero on the diagonal, so factorization without row
exchanges fails on essentially every netlist. If no candidate exists the
trailing matrix has an empty column and the matrix is singular; the failing
column index is returned and the error message names the node or branch,
which turns the most common user error, a node with no DC path, into a
self-explaining diagnostic.

What is missing, deliberately, is fill-reducing ordering. Real SPICE uses
Markowitz ordering; here the elimination order is netlist order, and the
worst case degrades toward dense cost. For tree-like and ladder-like
circuits, which is what a 1000-line simulator will realistically see, the
netlist order is close to optimal anyway: a 2000-node RC ladder solves its
operating point plus an AC sweep in tens of milliseconds.

\section{Newton-Raphson and the operating point}
\loc{device models, plus the solver drivers}{50}{572}

Nonlinear devices enter the linear system through companion linearization.
A device current $i(v)$ is replaced at the current iterate $v^{(k)}$ by
\[
i \approx g\, v + I_{eq}, \qquad
g = \left.\frac{\partial i}{\partial v}\right|_{v^{(k)}}, \qquad
I_{eq} = i(v^{(k)}) - g\, v^{(k)},
\]
that is, a conductance in parallel with a current source. Assembling and
solving this linearized system is one Newton iteration. Convergence is
declared when every unknown moves less than
$\text{tol}_i = \tau_i + \text{RELTOL} \cdot \max(|x_i^{(k)}|, |x_i^{(k+1)}|)$
between iterates, with $\tau_i = \text{VNTOL} = 1\,\mu\text{V}$ for node
voltages and $\text{ABSTOL} = 1\,\text{pA}$ for branch currents.

\paragraph{Junction limiting.}
The diode current $i = I_S (e^{v/V_T} - 1)$ makes plain Newton diverge: an
early iterate easily proposes $v = 5\,$V, and $e^{5/0.02585}$ overflows any
linearization. The classic fix is \texttt{pnjlim}, which damps the proposed
junction voltage logarithmically once it passes the critical voltage
\[
v_{crit} = V_T \ln \frac{V_T}{\sqrt{2}\, I_S},
\]
the point of maximum curvature of the exponential characteristic:

\begin{lstlisting}
fn pnjlim(vnew: f64, vold: f64, vt: f64, vcrit: f64) -> f64 {
    if vnew > vcrit && (vnew - vold).abs() > 2.0 * vt {
        if vold > 0.0 {
            let arg = 1.0 + (vnew - vold) / vt;
            if arg > 0.0 { vold + vt * arg.ln() } else { vcrit }
        } else {
            vt * (vnew / vt).ln()
        }
    } else {
        vnew
    }
}
\end{lstlisting}

\paragraph{Limiting must deny convergence.}
This is the one place where nanospice shipped a real bug worth documenting.
With limiting alone, the following happens on a common-emitter amplifier:
the voltage divider puts the base at 2.1\,V while the emitter is at ground,
\texttt{pnjlim} clamps the base-emitter voltage seen by the stamps to
0.11\,V, the transistor therefore conducts nothing, the node voltages stop
changing, and the delta-based convergence test reports success with the
transistor completely off. The solution is the classic one: any iteration
in which a limiter actually changed a junction voltage is not allowed to
converge. The assembly reports a \texttt{limited} flag and the Newton loop
requires it to be false. With that in place the same amplifier walks the
junction up in logarithmic steps and lands on the correct bias in a dozen
iterations.

\paragraph{Fallbacks.}
When plain Newton fails, two continuation methods run in order. Gmin
stepping solves the circuit with a large conductance $g_{min}$ across every
junction and reduces it decade by decade from $10^{-2}$ to $10^{-12}$\,S,
warm-starting each solve from the last. Source stepping scales all source
values by a factor ramped from 0.1 to 1.0; the zero-source circuit is
trivially solvable and each step tracks the solution continuously. Gmin
stepping handles floating-junction trouble, source stepping handles
feedback and bistable topologies, and together they are the standard SPICE
fallback chain:

\begin{lstlisting}
// second fallback: source stepping, ramp all sources up with warm starts
reset(x, lim);
for k in 1..=10 {
    let fac = k as f64 / 10.0;
    if newton(ckt, x, lim, &Mode::Dc { fac }, GMIN, ITL_OP).is_none() {
        die("operating point did not converge");
    }
}
\end{lstlisting}

\section{Device models}
\loc{device stamps, inside the solver section}{179}{751}

\paragraph{Diode.}
Shockley equation with emission coefficient $n$:
$i = I_S (e^{v/nV_T} - 1)$, conductance $g_d = I_S e^{v/nV_T} / nV_T$, plus
$g_{min}$ in parallel. The junction voltage passes through \texttt{pnjlim}
with per-device memory of the last limited value.

\paragraph{MOSFET, level 1.}
The square-law model with channel-length modulation:
\[
i_D = \begin{cases}
0 & v_{GS} \le V_{T0} \\
K_P \left( v_{ov} v_{DS} - \tfrac{1}{2} v_{DS}^2 \right)(1 + \lambda v_{DS})
  & v_{DS} < v_{ov} \\
\tfrac{1}{2} K_P v_{ov}^2 (1 + \lambda v_{DS}) & v_{DS} \ge v_{ov}
\end{cases}
\qquad v_{ov} = v_{GS} - V_{T0}.
\]
The model is evaluated in an \emph{effective frame} that handles both the
pmos polarity and the source-drain swap for $v_{DS} < 0$. With the sign
$s = \pm 1$ for nmos/pmos and effective terminals $(d_e, s_e)$ chosen so
that $v_{DS}^{e} \ge 0$, the actual terminal current is $I = s\, i_D$ and
the chain rule gives, for example,
\[
\frac{\partial I}{\partial v_g}
 = s \cdot \frac{\partial i_D}{\partial v_{GS}^{e}}
     \cdot \frac{\partial v_{GS}^{e}}{\partial v_g}
 = s \cdot g_m \cdot s = g_m .
\]
The two signs cancel in every derivative, so the conductance stamps are
identical for all four cases and only the right-hand side carries $s$. That
one observation removes what is otherwise a four-way case split:

\begin{lstlisting}
fn mos_op(dv: &Dev, x: &[f64]) -> (usize, usize, f64, f64, f64) {
    let Dev::M { d, g, s, kp, vt0, lambda, pmos } = dv else { unreachable!() };
    let sg = if *pmos { -1.0 } else { 1.0 };
    let (ed, es) = if sg * (x[*d] - x[*s]) >= 0.0 { (*d, *s) } else { (*s, *d) };
    let vgs = sg * (x[*g] - x[es]);
    let vds = sg * (x[ed] - x[es]);
    let (id, gm, gds) = mos_eval(vgs, vds, *kp, *vt0, *lambda);
    (ed, es, sg * id, gm, gds)
}
\end{lstlisting}

\paragraph{BJT.}
The Ebers-Moll model in transport form. With
$e_f = e^{v_{BE}/V_T}$ and $e_r = e^{v_{BC}/V_T}$:
\[
i_{CT} = I_S (e_f - e_r), \qquad
i_{BE} = \frac{I_S}{\beta_F}(e_f - 1), \qquad
i_{BC} = \frac{I_S}{\beta_R}(e_r - 1),
\]
\[
I_C = i_{CT} - i_{BC}, \qquad
I_B = i_{BE} + i_{BC}, \qquad
I_E = -(I_C + I_B).
\]
Four small-signal conductances
$g_{mf} = \partial i_{CT}/\partial v_{BE}$,
$g_{mr} = -\partial i_{CT}/\partial v_{BC}$,
$g_\pi = \partial i_{BE}/\partial v_{BE}$,
$g_\mu = \partial i_{BC}/\partial v_{BC}$
fill the $3 \times 3$ Jacobian stamp
\[
\begin{pmatrix}
\partial I_C / \partial v_b & \partial I_C / \partial v_c & \partial I_C / \partial v_e \\
\partial I_B / \partial v_b & \partial I_B / \partial v_c & \partial I_B / \partial v_e \\
\partial I_E / \partial v_b & \partial I_E / \partial v_c & \partial I_E / \partial v_e
\end{pmatrix}
=
\begin{pmatrix}
g_{mf} - g_{mr} - g_\mu & g_{mr} + g_\mu & -g_{mf} \\
g_\pi + g_\mu & -g_\mu & -g_\pi \\
g_{mr} - g_{mf} - g_\pi & -g_{mr} & g_{mf} + g_\pi
\end{pmatrix},
\]
whose rows and columns each sum to zero, a quick correctness check for any
MNA stamp. The pnp uses the same npn-frame equations through the sign trick
above; both junctions carry their own \texttt{pnjlim} memory, which is why
the limiter state is two values per device. An integration test builds the
same common-emitter stage as npn and as pnp mirror image and requires the
solutions to be exactly symmetric.

\section{Transient analysis}
\loc{analyses, with output selection}{217}{999}

\paragraph{Companion models.}
The trapezoidal rule applied to the capacitor equation $i = C\,dv/dt$ over a
step $h$ gives
\[
\frac{i_{n+1} + i_n}{2} = C\, \frac{v_{n+1} - v_n}{h}
\quad\Longrightarrow\quad
i_{n+1} = \underbrace{\frac{2C}{h}}_{g_{eq}} v_{n+1}
        - \underbrace{\left( \frac{2C}{h} v_n + i_n \right)}_{-I_{eq}},
\]
a conductance $g_{eq}$ in parallel with a history current source: the
capacitor becomes a resistor for the duration of one step. The inductor is
dual, $v_{n+1} = (2L/h)(i_{n+1} - i_n) - v_n$, stamped into its branch row.
Each reactive device keeps a two-entry state, its value and the dual
quantity, updated after every accepted step. The trapezoidal rule is
A-stable and undamped; an integration test runs an LC tank for twenty
periods and requires the amplitude to survive.

\paragraph{Timestep control.}
The local truncation error of the trapezoidal rule is
$\varepsilon_{tr} = -\tfrac{h^3}{12} \dddot{x}(\xi)$. To estimate
$\dddot{x}$ without extra solves, the integrator carries a quadratic
polynomial predictor through the last three accepted points. With step
history $d_1, d_2$ and the new step $h$, Lagrange extrapolation gives
coefficients
\[
\ell_0 = \frac{(h + d_1)(h + d_1 + d_2)}{d_1 (d_1 + d_2)}, \quad
\ell_1 = -\frac{h (h + d_1 + d_2)}{d_1 d_2}, \quad
\ell_2 = \frac{h (h + d_1)}{(d_1 + d_2)\, d_2},
\]
and the predictor error is
$\varepsilon_p = \tfrac{1}{6} \dddot{x} \, h (h + d_1)(h + d_1 + d_2)$.
Corrector and predictor bracket the same third derivative, so their
difference measures it (a Milne-style estimate), and the trapezoidal share
of the gap is
\[
\varepsilon_{tr} \approx
\frac{h^3/12}{\;h(h+d_1)(h+d_1+d_2)/6 + h^3/12\;} \;\bigl( x_c - x_p \bigr),
\]
which for uniform steps reduces to the classic factor $1/13$. The error is
compared per unknown against $\text{TRTOL} \cdot \text{tol}_i$ with
$\text{TRTOL} = 7$, the step is rejected above 1, and the next step follows
the cube-root rule $h' = 0.9\, h\, r^{-1/3}$, clamped to $[0.3, 2] \times h$
and capped at \texttt{tstep}:

\begin{lstlisting}
let etrap = hs * hs * hs / 12.0;
let fac = etrap / (hs * (hs + d1) * (hs + d1 + d2) / 6.0 + etrap);
let mut r: f64 = 0.0;
for i in 1..m {
    r = r.max(fac * (xn[i] - xp[i]).abs() / (TRTOL * tolv(i, nn, xn[i], x[i])));
}
let scale = 0.9 / r.max(1e-8).cbrt();
\end{lstlisting}

The predictor doubles as the Newton starting point, which is free and
usually saves an iteration. Before three points of history exist, and after
a nonconvergent step (cut by 8, the SPICE2 reflex), the iteration count
controls the step instead.

\paragraph{Breakpoints.}
A pulse edge in the middle of a step is invisible to the method: the
companion history current at the interval start belongs to the pre-edge
circuit, and the predictor extrapolates smoothness that is not there. The
first RC test against the analytic step response failed at exactly the
first output point for this reason. The classic remedy is exact: every
corner of every pulse source (delay, rise end, fall start, fall end, per
period) is registered as a breakpoint, steps are clamped to land on
breakpoints exactly, and after crossing one the integrator restarts with a
small step and cleared predictor history, since the waveform derivative is
discontinuous there.

\section{Small-signal AC}
\loc{numbers: Num trait and complex}{78}{101}

AC analysis linearizes the circuit at the operating point and solves
$(G + j\omega C)\, \tilde{x} = b$ per frequency, where $b$ holds the
\texttt{ac} source values. The implementation leans on one observation: at
a converged operating point the limiters pass voltages through unchanged, so
\emph{the resistive matrix $G$ is exactly the DC Jacobian that the last
Newton iteration already assembles}. The AC path therefore reuses
\texttt{assemble} verbatim, maps the result to complex, and adds only the
reactive stamps $j\omega C$ and the branch entry $-j\omega L$. No device
model contains AC-specific code.

The complex solve is the same 96-line sparse LU. The solver is generic over
a three-method trait, implemented for \texttt{f64} and a 15-line complex
type:

\begin{lstlisting}
trait Num:
    Copy
    + std::ops::Add<Output = Self>
    + std::ops::Sub<Output = Self>
    + std::ops::Mul<Output = Self>
    + std::ops::Div<Output = Self>
{
    fn zero() -> Self;
    fn one() -> Self;
    fn mag(self) -> f64;
}
\end{lstlisting}

\texttt{mag} is the pivoting magnitude, absolute value for reals and modulus
for complex. This trait plus operator implementations is the entire price
Rust charges for not shipping a complex type in std; the alternative, two
copies of the factorization, costs more lines and can drift apart.

\section{The parser}
\loc{parser}{262}{426}

The parser is the largest section and the least glamorous. Line discipline
follows SPICE: the first line is a title, \texttt{*} comments, \texttt{;}
trails, \texttt{+} continues, everything is case-insensitive, and
parentheses and commas are replaced by whitespace once during preprocessing,
which lets \texttt{sin(0 1 1k)} and \texttt{sin 0 1 1k} tokenize
identically.

Numbers use the SPICE grammar: the longest prefix that parses as a float,
then a unit suffix. The only subtlety is that \texttt{meg} must match before
\texttt{m}, handled by table order:

\begin{lstlisting}
fn num(tok: &str) -> Option<f64> {
    let t = tok.trim();
    if !t.is_ascii() {
        return None;
    }
    let mut end = t.len();
    while end > 0 && t[..end].parse::<f64>().is_err() {
        end -= 1;
    }
    if end == 0 {
        return None;
    }
    let v: f64 = t[..end].parse().unwrap();
    let mult = SUFFIXES.iter().find(|(s, _)| t[end..].starts_with(s)).map_or(1.0, |(_, m)| *m);
    Some(v * mult)
}
\end{lstlisting}

The ASCII guard exists because byte-indexed slicing panics inside multibyte
characters; a netlist containing a literal micro sign must produce
\emph{bad number}, not a crash, and a fuzz test holds the parser to that.

\texttt{.model} cards are collected in a prepass so they may appear anywhere
in the deck. A device line referencing a model has the model's parameter
tokens spliced in \emph{behind} its inline tokens; since parameter lookup
returns the first match, instance parameters override model parameters by
construction, with no explicit precedence logic.

Subcircuits are the notable omission. \texttt{.subckt} costs roughly a
hundred lines of name spacing and flattening and is the single feature most
requested of any SPICE clone; it lost to the transistor models.

\section{Validation}

The test suite is 18 integration tests running the built binary on real
netlists. Tests are free of the budget, so robustness lives there, and the
suite leans on references that are independent of the implementation:

\begin{itemize}
\item Analytic: RC and RL step responses against $1 - e^{-t/\tau}$, the RC
corner at $1/\sqrt{2}$ and $-45^\circ$, LC amplitude $i_0 \sqrt{L/C}$, and
energy conservation over twenty periods.
\item Bias points: MOSFET saturation current, BJT $I_C = \beta_F I_B$
consistency, diode KCL closure, and exact npn/pnp mirror symmetry.
\item Structural: a 200-node ladder and a 30-stage LCG-randomized resistor
ladder verified node by node against a Thevenin reduction computed
independently in the test.
\item Negative: seven malformed netlists (unknown devices, bad numbers,
missing tokens, non-ASCII input, floating nodes) must exit with a clean
named error and must not panic; the floating-node case must name the node.
\item Meta: the LOC counter itself.
\end{itemize}

The LTE controller has a dedicated regression: an RC circuit whose time
constant is ten times smaller than \texttt{tstep}, which a fixed-step
integrator fails badly and the controller must resolve on its own.

\section{The cut list}
\label{sec:cuts}

\begin{table}[h]
\centering
\begin{tabular}{llr}
\toprule
feature & why it lost & est. lines \\
\midrule
\texttt{.subckt} hierarchy & lost to the transistor models & $\sim$100 \\
Markowitz ordering & netlist order suffices at this scale & $\sim$80 \\
Gummel-Poon BJT & Ebers-Moll covers the teaching core & $\sim$80 \\
junction capacitances & needs charge state on every junction & $\sim$60 \\
noise analysis & needs the adjoint system & $\sim$80 \\
\texttt{.param} expressions & a calculator is its own project & $\sim$60 \\
current-based convergence check & limiting denial covers the failure mode & $\sim$15 \\
\bottomrule
\end{tabular}
\caption{Features that did not fit, with the reason and the estimated cost.}
\end{table}

On language: the device enum with exhaustive matching carries the parser,
the assembly, and the output selection, and is the main reason the budget
holds. C++ would trade the roughly 60 lines the \texttt{Num} trait and
complex type cost against visitor or inheritance boilerplate around
\texttt{std::variant}, roughly a wash with fewer static guarantees. Fortran
would shine on dense arrays, which this simulator does not use, and struggle
badly on the parser and the sparse row structure.

\section{Closing}

The final count is 999 of 1000 lines, with the classic algorithm set
complete: MNA, sparse LU with partial pivoting, Newton with pnjlim and both
continuation fallbacks, trapezoidal integration with LTE control and
breakpoints, small-signal AC, nine device types, four analyses, model
cards, initial conditions, and output selection. The one line of headroom
is left as a reminder that the next feature must pay for itself.

\end{document}
